% \section{Phase-space description of clustering}\label{sec:phase-space}
\section{Clarifying the initial conditions of the pushforward:
phase-space perspective}\label{sec:phase-space}
\noindent Since LPT is based on a
second-order ODE (Newton's equation),
phase space is the natural arena in which to describe the trajectories.
A key aspect of phase-space dynamics is that they are
determined by state, the initial positions $\q$ and initial velocities $\u$.
% In turn, the number density~\eqref{eq:n} is
% determined by counting up these trajectories.
This raises a question: why in Eq.~\eqref{eq:n} are the
trajectories initialized 
by positions only, and not both positions {and} velocities?
In other words, why is it $\x_t(\q)$ and not $\x_t(\q,\u)$?
Here we will clarify the conditions under which
Eq.~\eqref{eq:n} is valid, pointing out the implicit assumption
of cold initial conditions.
%We will find that there is an implicit
%assumption about the initial conditions being \emph{cold}.

% In this section we show why positions are sufficient,
% and in doing so we will clarify the regime of validity of Eq.~\eqref{eq:n}.
% Here we enlarge the usual description to phase space,
% showing how kinematic information enters the spatial correlations.
% Our treatment also provides a new handle on
% modelling the correlation function without losing
% multistreaming effects.

\subsection{Pushforward from arbitrary initial conditions}
\noindent Begin with the phase-space density $f(\x,\v,t)$. This
gives the probability $f(\x,\v,t)\,\dif^3\x\,\dif^3\v$
of finding a particle in a phase-space volume $\dif^3\x\,\dif^3\v$ 
centred at $(\x,\v)$.
The density is $\rho(\x,t)=\int\dif^3\v\, f(\x,\v,t)$.
By the chain rule it is convenient to write
\be\label{eq:f-chain}
f(\x,\v,t)=\rho(\x,t)\sp p(\v\Mid\x,t)\,,
\ee
where $p(\v\Mid\x,t)$ is the velocity distribution
at position $\x$ and time $t$.
% by assuming that we have
% complete information about the system, as given by
% the Klimontovich distribution function
% \be\label{eq:fK}
% f(\x,\v,t)
% =\sum_{a}\delD[\x-\x_a(t)]\,\delD[\v-\v_a(t)]\,,
% \ee
% Here $\v_a(t)=a(t)\dot\x_a(t)$ is the $a$th particle (peculiar) velocity
% (this is equal to the canonical momenta with unit mass).
% Note that in the limit $N\to\infty$, the discrete 
% distribution~\eqref{eq:fK} can be replaced by
% the smooth phase-space distribution (as in the Vlasov equation).
% For convenience we define the phase-space position $\w=(\x,\v)$.
% The phase-space trajectory 
% \be
% \w_a(t)=(\x_a(t),\v_a(t))
% \ee
% of the
% $a$th particle  is governed by a system of first-order ODEs
% (e.g.\ Hamilton's equation) of the form
% ${\dif\w}/{\dif t}=\mbf{F}(\w,t)$, $\w=(\x,\v)$, whose solution we
% denote $\bm\Phi_t(\w_0)$  with initial condition $\w_0=(\x_0,\v_0)$.
% The important step in what follows is the lifting of the
% trajectories $\bm\Phi_t$ into $f(\w,t)\equiv f(\x,\v,t)$
% by writing 
Similar to Eq.~\eqref{eq:n}, the phase-space density can be expressed in
terms of the initial density $f_0(\q,\u)\equiv f(\q,\u,0)$ as
% [with $f(\X,t)=f(\x,\p,t)$]
\be\nonumber %\label{eq:f}
f(\x,\v,t)
% =\int_{\w_0}\delD\big[\w-\bm\Phi_t(\w_0)\big]\ssp f_0(\w_0)\,,
=\int\dif^3{\q}\int\dif^3{\u}\;
\delD\big[\x-\x_t(\q,\u)\big]\ssp
\delD\big[\v-\v_t(\q,\u)\big]\ssp f_0(\q,\u)\,,
\ee
where $\x_t$ and $\v_t$ solve Hamilton's equation.
% This is equivalent to the above standard form~\eqref{eq:fK}, as is easily
% checked by substituting $f_0(\w_0)=\sum_a\delD(\w_0-\w_{a,0})$ and
% evaluating the integral over $\w_0$.
This expression can thus be seen as the formal solution to
the Vlasov--Poisson equation.
Note that due to self-consistency $\x_t(\q,\u)$ and $\v_t(\q,\u)$
also depend on the initial density $f_0$; for brevity
this dependence has been suppressed.
% Ensemble averages of any observable of $f(\x,\v,t)$
% is taken over realizations of the
% initial three-dimensional phase-space sheet, or equivalently
% over $n_0$ (or $\delta_0$)~\citep{Abel:2011ui,Shandarin:2012}.
% \footnote{
% The initial position and velocity of each particle $\q$
% is specified by the map~\citep{Abel:2011ui,Shandarin:2012}
% \be
% \q\mapsto (\x_0,\u_0)=(\q+D_0\ssp\bm\psi_0(\q),\dot{D}_0\ssp\bm\psi_0)\,,
% \ee
% where $D_0=D(0)$ is the growth factor of the linear growing mode,
% and $\bm\psi_0(\q)$ is the initial displacement field.
% This map describes a smooth three-dimensional sheet in phase space.
% The initial density and velocity fields are
% related to the displacement field $\bm\psi_0(\q)$ by
% $\delta_0=-D(0)\bm\nabla\cdot\bm\psi_0$ and
% $\mbf{U}_0\equiv \dot{D}(0)\ssp\bm\psi_0$.
% }



%%%%%%%%%%%%%%%%%%%%%%%%%%%%
% The notation above hides that
% Eq.~\eqref{eq:f} is a
% \emph{nonlinear} functional of $f_0$.
% Since the acceleration of a particle is determined by interactions with the rest of the system (through Poisson's equation),
% $\bm\Phi_t$ implicitly depends on the state of the entire 
% distribution (specifically, the positions of all other particles). As a self-consistent model, the initial position and velocity $\w_0$
% of any given particle are not enough to determine its future trajectory.
% (If $\mbf{F}$ is a fixed vector field then $\w_0$ would be
% sufficient, but this would not be a self-consistent theory.)
% But knowing the initial distribution $f_0$
% is enough to fix the trajectories of all particles, and so
% we might have written instead $\bm\Phi_t[f_0](\w_0)$. 
% Another way to understand this is because the transport
% equation $\partial_t f+\bm\nabla\cdot(\mbf{F}f)=0$ is nonlinear
% in $f$, given that the
% force $\mbf{F}=\mbf{F}[f]$ by Poisson's equation.
% However, even though solving this integro-differential
% equations is perhaps only possible numerically,
% we can still formally write the
% general solution for $f(\w,t)$ as Eq.~\eqref{eq:f}.

% This is still progress; Eq.~\eqref{eq:f} is
% an explicit expression in the initial conditions (appearing
% only on the right-hand side) which is ultimately what
% we want to compute expectations over.
% And even though we do not have explicit solutions for the
% trajectories $\bm\Phi_t$, this is not a problem; we
% care only about the statistics of $\bm\Phi_t$ rather than the detailed solutions for any given realization of the initial conditions.
%%%%%%%%%%%%%%%%%%%%%%%%%%%%

% \bea\label{eq:n-lifted}
% n(\x,t)
% % &=\int_\v\int_{\w_0}
% %     \delD[\w-\bm\Phi_t(\w_0)]\ssp f_0(\w_0) \\
% % &=\int_\v\int_{\w_0}
% %     \delD[\x-\x_t(\w_0)]\ssp
% %     \delD[\v-\v_t(\w_0)]\ssp f_0(\w_0)     \nonumber\\
% =\int_\v f(\x,\v,t)
% &=\int_\q\int_\u
%     \delD[\x-\x_t(\q,\u)]\ssp f_0(\q,\u)\,,
% \eea
% where the second equality follows from inserting
% Eq.~\eqref{eq:f}. Since by the chain rule the distribution
% function can at any time be factorized as
% $f(\x,\v,t)=n(\x,t)\ssp p(\v\Mid\x,t)$,
% Eq.~\eqref{eq:n-lifted} can be written
Now inserting the foregoing equation into
$\rho(\x,t)=\int\dif^3\v\,f(\x,\v,t)$
and using Eq.~\eqref{eq:f-chain} to write
$f_0(\q,\u)=\rho_0(\q)\ssp p_0(\u\Mid\q)$, we have
\be\label{eq:n2}
\rho(\x,t)=\int\dif^3\q\; T_t(\x\Mid\q)\, \rho_0(\q)\,,
% n(\x,t)=\int_\q\bigg(\int_\u p_0(\u\Mid\q)\ssp
%     \delD\big[\x-\x_t(\q,\u)\big]\bigg)\, n_0(\q)\,.
\ee
with the transition kernel
\be\label{eq:T_t}
T_t(\x\Mid\q)
\equiv\int\dif^3\u\, p_0(\u\Mid\q)\,
    \delD\big[\x-\x_t(\q,\u)\big]\,,
\ee
which is normalized: $\int\dif^3\x\, T_t(\x\Mid\q)=1$.
Equation~\eqref{eq:n2} is the generalization of Eq.~\eqref{eq:n}
to arbitrary initial conditions at any time.
If particles are multistreaming at $\q$,
then kinematic information is lost in the velocity-averaged
kernel~\eqref{eq:T_t}.
In this case tracking positions only is not sufficient to determine
the future trajectory.
%\footnote{The average over $\u$ should {not} be mistaken for
%having the effect of a smoothing kernel on the sharp transitions.
%Transitions remain sharp (before ensemble averaging).
%This is because of phase mixing, whereby
%the initial phase-space sheet repeatedly folds in on itself.
%Although this leads to a distribution of velocities
%at a given position, the number of folds is not continuous number,
%meaning that $p_0(\u\Mid\q)$ and thus $T_t(\x\Mid\q)$
%is generally a spiky distribution, not a smooth one.}

%This kernel depends on the initial conditions.
% As a result $T_t(\x\Mid\q)$ is a multiply spiked
% distribution (multiple $\x$ possible for a given $\q$).
% For example, suppose there are two streams with velocity
% $\u_1$ and $\u_2$ at $\q$:
% \be
% p_0(\u\Mid\q)= \tfrac12[\delD(\u-\u_1)+\delD(\u-\u_2)]
% \ee
% Plugging this into Eq.~\eqref{eq:T_t} we see that a particle at $\q$
% can land on one of two Eulerian positions, depending on the
% initial velocity.
% Then
% \be
% T_t(\x,\q)=\tfrac12\delD[\x-\x_t(\q,\u_1)]
%     +\tfrac12\delD[\x-\x_t(\q,\u_2)]
% \ee
% The integration over integral generally broadens the sharp transition kernel
% and has the effect of smoothing out the density peaks.


%The transition probability of a pair $(\q_1,\q_2)$ at
%an arbitrary initial time is generally given by
%\be\label{eq:pt-v-avg}
%p_t(\x_1,\x_2\Mid\q_1,\q_2) 
%=\llangle T_t(\x_1\Mid\q_1)\sp T_t(\x_2\Mid\q_2)\rrangle\,,
%\ee
%where the density-weighted ensemble average is again over the
%initial density $n_0$
%(or $\delta_0$) on which $T_t(\x\Mid\q)$ implicitly depends through
%$\x_t$ and $p_0(\u\Mid\q)$. 
%The two-point function $\xi_n(\x_1,\x_2,t)$ is given as before
%by Eq.~\eqref{eq:MASTER}, but now with the
%transition probability~\eqref{eq:pt-v-avg}.



\subsection{Pushforward from cold initial conditions}
\noindent Equation~\eqref{eq:n} is recovered from
the general formula~\eqref{eq:n2} in the case of cold initial conditions.
Such conditions require that the pushforward be applied at an early
epoch when the Universe was well approximated by a pressureless
perfect fluid, described
by the distribution function
$f_0(\q,\u)=\rho_0(\q)\ssp\delD[\u-\v_{m0}(\q)]$ with
the initial velocity field $\v_{m0}(\q)\propto\nabla\nabla^{-2}\delta_0$
corresponding to the growing mode.
Comparing this with Eq.~\eqref{eq:f-chain}, we deduce
$p_0(\u\Mid\q)=\delD[\u-\v_{m0}(\q)]$
which, when inserted into Eq.~\eqref{eq:T_t}, yields
$T_t(\x\Mid\q)=\delD[\x-\x_t(\q,\v_{m0}(\q)]$.
This transition kernel recovers Eq.~\eqref{eq:n} in a form
that makes clear that the trajectory also depends on
initial velocity too.
%Thus Eq.~\eqref{eq:n2} reads
%\be\label{eq:n-alt}
%\rho(\x,t)
%=\int_{\q}
%    \delD\big[\x-\x_t(\q,\v_{m0}(\q))\big]\, \rho_0(\q)
%\ee
%This expression is equivalent to Eq.~\eqref{eq:n}. now with
%the hidden dependence on the initial velocity $\u=\v_{m0}(\q)$
%made explicit.
In fact kinematic information is already contained in the
initial density $\rho_0$, which is conventionally suppressed anyway.
The phase-space picture makes this clear.
For cold initial conditions the initial distribution function $f_0$
is confined to a three-dimensional sheet, essentially $n_0$,
in six-dimensional phase space~\citep{Abel:2011ui,Shandarin:2012}.%
\footnote{More precisely,
for a given linear $\delta_0$ the three-dimensional sheet is described by the
map from configuration space into phase space:
\be
\q\mapsto (\x_0,\v_0)
=\big(\q+D_0\ssp\bm\psi_0(\q),\dot{D}_0\ssp\bm\psi_0(\q)\big)\,,
\ee
where $D_0=D(0)$ and $\bm\psi_0(\q)\propto\nabla\nabla^{-2}\delta_0$
is the initial displacement field.
%Note that by the continuity equation
%$\bm\psi_0(\q)\propto\nabla\nabla^{-2}\delta_0$.
%The initial density and velocity fields are
%related to $\bm\psi_0(\q)$ by $\delta_0=-D_0\nabla\cdot\bm\psi_0$
%and $\v_{m0}\equiv \dot{D}_0\ssp\bm\psi_0$.
}
Sampling this sheet yields the initial positions and
velocities for the dark matter particles in an $N$-body simulation.

The above considerations relate to a particle (or multistreaming)
description of clustering.
In standard perturbation theory~\citep{Bernardeau:2001qr},
the stronger assumption is made that cold dark matter is
a pressureless perfect fluid described by
$f(\x,\v,t)=\rho(\x,t)\,\delD[\v-\v_{m}(\x,t)]$ for all times.
In this case, $\x_t(\q)$ is given by the equation of motion
\be\label{eq:eom-ss}
\frac{\dif\x_t(\q)}{\dif t}=\frac{1}{a}\,\v_m(\x_t(\q),t)\,,
\ee
and the pushforward~\eqref{eq:n} can be applied at {any} time.
Although single streaming is not required, we show in
Appendix~\ref{app:path-integral} that the
pushforward can be written in path-integral form.

